% English translation of checked_pages/esq_119_007.tex; source: cote 119, site 7, PDF page 8, typescript page 5.
% AI-translated by GPT-6 Astra (OpenAI, ChatGPT Pro mode).
\scanpage{119/7}
\mspage{5}

of $\hat{\pi}_{0,3}$, i.e. on the corresponding parameter or parameters, for it
to come from an element of $\mathrm{I}\!\Gamma$ - which would provide a
``purely algebraic'' description, in terms of profinite groups and without
reference to the Galois theory of number fields, of the Galois
group $\mathrm{I}\!\Gamma = \mathrm{Gal}(\bar{\mathbb{Q}}/\mathbb{Q})$ !

Perhaps even a conjectural characterization of $\mathrm{I}\!\Gamma$
as a subgroup of Autext($\hat{\pi}_{0,3}$) is for the moment out of
reach($^1$) ; I have no conjecture to propose yet. Another
task, however, can be approached immediately: describing
the action of $\mathrm{I}\!\Gamma$ on the entire Teichmüller tower in terms of its
action on the ``first level'' $\hat{\pi}_{0,3}$ , i.e. expressing an automorphism
of this tower in terms of the ``parameter'' in $\hat{\pi}_{0,3}$ that specifies
the varying element $\gamma$ of $\mathrm{I}\!\Gamma$. This is connected with a representation of the
Teichmüller tower (as a groupoid equipped with a ``gluing''
operation) by generators and relations, which will in particular give
a presentation by generators and relations, in the ordinary sense, of
each of the $\hat{T}_{g,\nu}$ (as a profinite groupoid). Here, even for
$g = 0$ (thus when the corresponding Teichmüller groups are
``well-known'' braid groups), the generators and relations
known to date that have come to my attention seem unusable
as they stand, because they lack the invariance and
symmetry properties indispensable for the action of $\mathrm{I}\!\Gamma$ to be directly
readable from this presentation. This is connected in particular with the fact that
people still persist, when calculating with fundamental
groups, in fixing a single base point, instead of cleverly choosing
a whole collection invariant under the symmetries of the
situation, which are thereby lost along the way. In certain situations
(such as descent theorems in the manner of Van Kampen for fundamental
groups), it is much more elegant, even indispensable\frapparatus{V-FR-007-001} for understanding
anything about them, to work with fundamental groupoids
relative to a suitable collection of base points , and this is certainly
so for the Teichmüller tower. It would seem (incredible,
but true !) that the very geometry of the first level of the
Teichmüller tower (thus corresponding to ``moduli'' either for projective

\margin{7}

